% !TEX TS-program = xelatex
% !TEX encoding = UTF-8 Unicode
% -*- coding: UTF-8; -*-
% vim: set fenc=utf-8

%%%%%%%%%%%%%%%%%%%%%%%%%%%%%%%%%%%%%%%%%%%%%%%%%%%%%%%%%%%%%%%%%
%% CV.tex — Inseong Jung
%% 템플릿: https://github.com/zachscrivena/simple-resume-cv (Unlicense)
%%
%% 빌드: latexmk -xelatex CV.tex
%% (pdflatex 로는 컴파일되지 않는다. XeLaTeX 전용)
%%%%%%%%%%%%%%%%%%%%%%%%%%%%%%%%%%%%%%%%%%%%%%%%%%%%%%%%%%%%%%%%%

\documentclass[a4paper,MMMyyyy,nonstopmode]{simpleresumecv}
% 클래스 옵션: a4paper, letterpaper, nonstopmode, draftmode
%              MMMyyyy, ddMMMyyyy, MMMMyyyy, ddMMMMyyyy, yyyyMMdd, yyyyMM, yyyy

%%%%%%%%%%%%%%%%%%%%%%%%%%%%%%%%%%%%%%%%%%%%%%%%%%%%%%%%%%%%%%%%%
%% PREAMBLE
%%%%%%%%%%%%%%%%%%%%%%%%%%%%%%%%%%%%%%%%%%%%%%%%%%%%%%%%%%%%%%%%%

\newcommand{\CVAuthor}{Inseong Jung}
\newcommand{\CVTitle}{Inseong Jung — Curriculum Vitae}
\newcommand{\CVNote}{}
\newcommand{\CVWebpage}{https://github.com/jung-inseong}

%% 한글 지원.
%% 템플릿 기본 폰트(Tinos, GNU FreeFont)에는 한글 글리프가 없다.
%% ucharclasses 가 한글 문자를 만나면 자동으로 아래 폰트로 바꿔주므로
%% 본문에서 한글을 따로 감쌀 필요가 없다.
\usepackage{ucharclasses}
\newfontfamily\KoreanFont{Malgun Gothic}
\setTransitionsFor{HangulSyllables}{\KoreanFont}{\normalfont}
\setTransitionsFor{HangulCompatibilityJamo}{\KoreanFont}{\normalfont}

%% 마침표 뒤 여분 공백을 없앤다. 발표 목록처럼 마침표가 많은 항목에서
%% 줄이 불필요하게 늘어나 오른쪽 끝의 날짜가 다음 줄로 밀리는 것을 막는다.
\frenchspacing

%% PDF 속성.
\hypersetup{
pdftitle={\CVTitle},
pdfauthor={\CVAuthor},
pdfsubject={Curriculum Vitae},
pdfcreator={XeLaTeX},
pdfproducer={},
pdfkeywords={sociology, political polarization, computational sociology},
unicode=true,
bookmarks=true,
bookmarksopen=true,
pdfstartview=FitH,
pdfpagelayout=OneColumn,
pdfpagemode=UseOutlines,
hidelinks,
breaklinks}

%%%%%%%%%%%%%%%%%%%%%%%%%%%%%%%%%%%%%%%%%%%%%%%%%%%%%%%%%%%%%%%%%
%% DOCUMENT
%%%%%%%%%%%%%%%%%%%%%%%%%%%%%%%%%%%%%%%%%%%%%%%%%%%%%%%%%%%%%%%%%

\begin{document}

%%%%%%%%%%%%%%%
% TITLE BLOCK %
%%%%%%%%%%%%%%%

\Title{Inseong Jung 정인성}

\begin{SubTitle}
Department of Sociology, Korea University, Seoul, South Korea
\par
\href{mailto:prsnlity01@korea.ac.kr}
{prsnlity01@korea.ac.kr}
\end{SubTitle}

\begin{Body}

%%%%%%%%%%%%%%%
%% EDUCATION %%
%%%%%%%%%%%%%%%

\Section
{Education}
{Education}
{PDF:Education}

\Entry
\textbf{Korea University}, Seoul, South Korea

\Gap
\BulletItem
Master of Arts in Sociology (Expected)
\hfill
\DatestampY{2025} --
\DatestampY{2027}

\Gap
\BulletItem
Bachelor of Arts in Sociology and Bachelor of Economics
\hfill
\DatestampY{2020} --
\DatestampY{2025}
\begin{Detail}
\SubBulletItem
Accelerated B.A.--M.A. Program in Sociology
\SubBulletItem
Early Graduation with High Honors
\end{Detail}

%%%%%%%%%%%%%%%%%%%%%%
%% RESEARCH INTEREST %%
%%%%%%%%%%%%%%%%%%%%%%

\Section
{Research Interests}
{Research Interests}
{PDF:ResearchInterests}

\Entry
Cultural Schema, Political Polarization, Computational Sociology, Social Network

%%%%%%%%%%%%%%%%%%%%%%%%%%%%
%% CONFERENCE PRESENTATIONS %%
%%%%%%%%%%%%%%%%%%%%%%%%%%%%

\Section
{Conference Presentations}
{Conference Presentations}
{PDF:ConferencePresentations}

\Entry
\BulletItem
Jung, Inseong and Eun Kyong Shin.
``Have the US Partisans Become More Predictable? A Machine Learning Approach.''
2026 ASA Annual Meeting, Political Sociology Section Roundtable.
New York, United States.
\hfill
\DatestampYM{2026}{08}

\Gap
\BulletItem
Jung, Inseong.
``Do You Talk the Way You Think? Multidimensional Belief System and Unidimensional Talk.''
2026 Association of Korean American Sociologists Conference.
New York, United States.
\hfill
\DatestampYM{2026}{08}

\Gap
\BulletItem
Jung, Inseong.
``Multidimensionality of Political Polarization: A Belief System and Word Embedding Approach.''
Korean Sociological Association Conference.
Busan, South Korea.
\hfill
\DatestampYM{2026}{06}
\begin{Detail}
\SubBulletItem
정인성. ``정치적 양극화의 다차원성: 신념 체계와 발화 임베딩 연구.''
한국사회학회 전기사회학대회. 부산, 대한민국.
\end{Detail}

\Gap
\BulletItem
Jung, Inseong.
``How Do Attitudes Predict Partisan Identity? A Machine Learning Approach.''
Critical Sociological Association Conference.
Seoul, South Korea.
\hfill
\DatestampYM{2025}{11}
\begin{Detail}
\SubBulletItem
정인성. ``태도는 어떻게 당파성을 예측하는가? 머신러닝 활용 당파성 예측 연구.''
2025년 비판사회학대회. 서울, 대한민국.
\end{Detail}

\Gap
\BulletItem
Shin, Eun Kyong, Hannah Baek, and Inseong Jung.
``A Machine Learning Trilogy: Unraveling Adult Intelligence, Adolescent Academic
Achievement, and Political Identity.''
Harvard University Urban Data Workshop.
Massachusetts, United States.
\hfill
\DatestampYM{2025}{11}

\Gap
\BulletItem
Jung, Inseong.
``Martial Law and the Amplification of Affective Antagonism:
Online Community Posts Analysis for the Two Major Parties.''
Summer Conference of Critical Sociological Association of Korea.
Chungju, South Korea.
\hfill
\DatestampYM{2025}{08}
\begin{Detail}
\SubBulletItem
정인성. ``계엄과 정서적 적대감의 증폭: 양대 정당 온라인 커뮤니티 게시글 분석.''
비판사회학회 하계학술대회. 청주, 대한민국.
\end{Detail}

\Gap
\BulletItem
Jung, Inseong.
``Who Becomes an Adult? Effect of Regional Origin and Family Background
on Attaining Traditional Markers of Adulthood.''
Korean Sociological Association Conference.
Seoul, South Korea.
\hfill
\DatestampYM{2024}{12}
\begin{Detail}
\SubBulletItem
정인성. ``누가 어른이 되는가? 출신 지역과 가족 배경이 전통적 성인기 이행 지표
획득에 미치는 영향.'' 한국사회학회 후기사회학대회. 서울, 대한민국.
\end{Detail}

\Gap
\BulletItem
Jung, Inseong.
``Accident or State Violence? Social Disaster Discourse Analysis on
the Sewol Ferry Disaster and the Itaewon Tragedy.''
PNC 2024 Annual Conference and Joint Meetings.
Seoul, South Korea.
\hfill
\DatestampYM{2023}{08}

%%%%%%%%%%%%%%%%%%%%%%%%%%
%% GRANTS & SCHOLARSHIPS %%
%%%%%%%%%%%%%%%%%%%%%%%%%%

\Section
{Grants \& Scholarships}
{Grants \& Scholarships}
{PDF:GrantsScholarships}

\Entry
\BulletItem
Brain Korea 21 Graduate Scholarship, National Research Foundation
\hfill
\DatestampY{2026}

\Gap
\BulletItem
Student Researcher Project Grant, Rohmuhyun Foundation
\hfill
\DatestampY{2025}
\begin{Detail}
\SubBulletItem
Project: Cultural Infrastructure Disparity and Quality of Life
\SubBulletItem
KRW 3,000,000 (USD 2,200)
\SubBulletItem
노무현재단 2025년 학생주도 연구프로젝트. 문화기반시설 격차와 삶의 질.
\end{Detail}

\Gap
\BulletItem
Humanities 100-Year Scholarship, Korea Scholarship Foundation
\hfill
\DatestampY{2024}
\begin{Detail}
\SubBulletItem
Undergraduate Tuition Scholarship
\SubBulletItem
한국장학재단 인문 100년 장학금.
\end{Detail}

%%%%%%%%%%%%%%%%%%%%%%%%%%%
%% PROJECTS & EXPERIENCES %%
%%%%%%%%%%%%%%%%%%%%%%%%%%%

\Section
{Projects \& Experiences}
{Projects \& Experiences}
{PDF:ProjectsExperiences}

\Entry
\BulletItem
\textbf{Research Assistant}, Lab Sociomarkers, Korea University
\hfill
\DatestampY{2026} --
\begin{Detail}
\SubBulletItem
Principal Investigator: Eun Kyong Shin
\end{Detail}

\Gap
\BulletItem
\textbf{Student Researcher}, Korea University Alumni Association
\hfill
\DatestampY{2024} --
\DatestampY{2025}
\begin{Detail}
\SubBulletItem
Project: Korea University 120th Anniversary Alumni Engagement Survey
\SubBulletItem
Online survey design, data collection and analysis
\end{Detail}

\Gap
\BulletItem
\textbf{Research Assistant}, Korea Information Society Development Institute
\hfill
\DatestampY{2024}
\begin{Detail}
\SubBulletItem
Project: Governance for Responding to a Post-Truth Society
\SubBulletItem
Data research and organization
\end{Detail}

\Gap
\BulletItem
\textbf{Sergeant}, Korean Augmentation to the United States Army, ROK Army
\hfill
\DatestampY{2021} --
\DatestampY{2023}
\begin{Detail}
\SubBulletItem
RADAR Platoon, 2nd Battalion, 3rd FA Regiment, 1st ABCT, 1st Division, US Army
\SubBulletItem
Mandatory military service
\SubBulletItem
Cultural support and translation for the US armed forces in South Korea
\end{Detail}

\end{Body}

\end{document}
